% The proof, inspected -- drafted by loop0015 item 04 (2026-11-05). Every
% number of this section is printed, and every line citation re-checked
% against the read-only clone at commit 4eaee875, by
%   python3 -I paper3/scripts/proof_inspection.py
% which re-runs the two audit scripts of reports/scratch/eth_audit/
% (closure.py, defdiff.py) as subprocesses, performs the remaining checks
% itself, writes latex/tables/proof_inspection.tex, and exits non-zero if
% anything has moved. Unqualified Lean paths below are relative to
% lean/OAI/Algorithms/SubexponentialSAT/ in the clone; the challenge file is
% lean/ComparatorChallenges/SubexponentialSAT.lean. Nothing in the clone was
% built, imported or run for this section; the build is Section 5's.
\section{The proof, inspected}
\label{sec:proof}
\begingroup\sloppy

Section~\ref{sec:statement} read the theorem; this section reads the
proof, as far as reading can go. The question is not whether the
mathematics is right---the kernel decides that, and Section~\ref{sec:replay}
reports what the kernel said when we ran it---but whether anything in the
source could make the kernel's acceptance mean less than it appears to
mean: an axiom, a \texttt{sorry}, a trusted evaluation, an unsound
instance, a patched Mathlib, a statement that differs from the challenge's,
or a cost model of the proof's own that leaks into the theorem. We found
none of these. What we did find is that the repository offers no recorded
evidence that the closure type-checks at all---no \texttt{\#print axioms}
output, no continuous integration, no second kernel for this family---which
is why Section~\ref{sec:replay} exists. The last subsection joins the two.

\subsection{Method}
\label{sec:proof-method}

The clone of \texttt{openai/math} at commit \texttt{4eaee875} is treated as
untrusted data: every file was read with \texttt{cat}, \texttt{sed},
\texttt{grep} and \texttt{python3 -I}, and nothing in it was built,
imported or executed for this section. Two scripts from the audit of
2026-11-05 do the mechanical part: \texttt{closure.py} follows
\texttt{import} lines from the solution module, strips comments (nested
block comments included) and counts keyword occurrences in what remains;
\texttt{defdiff.py} normalises whitespace and compares each declaration of
the challenge file with the library's declaration of the same name. Both
were re-run for this section, and every number below is from that fresh
run; \repofile{paper3/scripts/proof_inspection.py} re-runs them, performs
the remaining checks, and regenerates Table~\ref{tab:proof-inspection}.
The inspection is of the spine of the proof---the final file, the
compiler, the entry lemmas, the constants---and of the whole closure by
keyword. The arithmetic beneath the spine was not read
(Section~\ref{sec:proof-not}).

\subsection{The closure}
\label{sec:proof-closure}

The comparator's JSON names one solution module,
\path{OAI.Algorithms.SubexponentialSAT.Main}, a sixteen-line aggregator
that imports \texttt{SqrtLog.ExecutionNativeFinal} and
\texttt{SqrtLogSpace.Final} and declares nothing
(\clonefile{Main.lean}). Its import closure inside the repository is
\textbf{3,091 modules}, every one with the prefix
\path{OAI.Algorithms.SubexponentialSAT}, totalling \textbf{550,551} raw
lines; the largest file is 2,873 lines
(\clonefile{SqrtLogSpace/DisjointReuse/SparsePriorityProgram.lean}). The
closure stops at \textbf{50} Mathlib modules and one module of Lean's own
library, \path{Lean.Elab.Tactic.Omega}, imported by
\clonefile{SqrtLogSpace/FiniteFallback.lean} (line~14) and
\clonefile{SqrtLogSpace/DisjointReuse/SparseCharging.lean} (line~5); it
is a tactic module, and the closure defines no tactic, macro or syntax of
its own (Section~\ref{sec:proof-keywords}). Nothing is imported from any
other family of the collection or from any of the repository's other
dependencies.

The family's directory holds \textbf{3,977} \texttt{.lean} files, so
\textbf{886} are outside the closure. They are not reachable from the
collection's root module either: \clonefile{lean/OAI.lean} imports the
family only at its line~49, \texttt{import} \path{OAI.Algorithms.SubexponentialSAT.Main}, and the lakefile's default target
is that root (\clonefile{lean/lakefile.lean}, line~89,
\texttt{@[default\_target] lean\_lib OAI}). They are therefore never built
by the default target and bear on nothing: dead files, not a hole. For
completeness, a raw search of the 886 for \texttt{sorry}, \texttt{axiom},
\texttt{native\_decide} and \texttt{admit} finds nothing, with or without
comments.

\subsection{The keyword sweep}
\label{sec:proof-keywords}

Table~\ref{tab:proof-inspection} gives the counts over the 3,091 files
with comments stripped. The ones that matter are the zeros. There is no
\texttt{axiom}, no \texttt{sorry}, no \texttt{admit}; no
\texttt{native\_decide} and no \texttt{ofReduceBool}, so nothing is
proved by trusting the compiled evaluator; no \texttt{implemented\_by},
\texttt{extern}, \texttt{unsafe} or \texttt{partial}, so no definition
has a hidden implementation and every recursion is one the kernel checks
(nine \texttt{termination\_by} and five \texttt{decreasing\_by} clauses);
no \texttt{elab}, \texttt{macro}, \texttt{syntax} or \texttt{simproc}, so
there is no metaprogramming that could rewrite goals outside the kernel's
view. The only \texttt{set\_option}s are two
\texttt{set\_option quotPrecheck false}
(\clonefile{SqrtLog/LeafProgramPreparedRecoveryQuery.lean}, line~15, and
\clonefile{SqrtLog/LeafProgramPreparedRecoveryWitness.lean}, line~14),
which switch off a pre-check on local notation and have no bearing on
soundness; there is no \texttt{maxHeartbeats}, \texttt{maxRecDepth} or
\texttt{autoImplicit} option anywhere in the closure, and the lakefile
sets \texttt{autoImplicit} to \texttt{false} for the whole package
(\clonefile{lean/lakefile.lean}, line~9).

Nine declarations are \texttt{opaque}, and every one carries a body:
\begin{itemize}
\item three recovery functions: \texttt{queryRecovery} at line~59 of
  \clonefile{SqrtLog/LeafProgramPreparedRecoveryQuery.lean}, with its
  \texttt{:= by} at line~86; \texttt{rawRecovery} and
  \texttt{fixedRecovery} at lines 178 and 198 of
  \clonefile{SqrtLog/LeafProgramPreparedRecoveryWitness.lean};
\item six \texttt{private opaque} constants of a subtype
  \texttt{\{ports // ports = \ldots\}}, each with an explicit value:
  line~22 of \clonefile{SqrtLog/LeafProgramQueryCleanup.lean} and five
  others.
\end{itemize}
An \texttt{opaque} with a body is a definition whose unfolding
is hidden from elaboration; the kernel still checks the body against the
type, so nothing is assumed. The 1,048 \texttt{noncomputable}
declarations, the eight \texttt{open Classical}, and the four explicit
uses each of \texttt{Classical.choice} and \texttt{propext} are the
classical reasoning the comparator permits. The 148 \texttt{instance}s
and 40 \texttt{local instance}s are \texttt{Fintype},
\texttt{Decidable} and algebraic instances; an instance cannot be
``inconsistent'' without an axiom, and there is none. The 14,034 lines
with \texttt{decide} are kernel-evaluated decision procedures, not the
native evaluator.

\begin{table}[t]
\centering\small
\caption{The inspection in numbers. Comments stripped where a keyword is
counted. Regenerate: \texttt{python3 -I paper3/scripts/proof\_inspection.py}.}
\label{tab:proof-inspection}
% Generated by python3 -I paper3/scripts/proof_inspection.py -- do not edit.
\begin{tabular}{@{}p{0.50\linewidth}p{0.17\linewidth}p{0.25\linewidth}@{}}
\toprule
quantity & value & script \\
\midrule
modules in the closure of \texttt{Main.lean} & 3,091 & \texttt{closure.py} \\
raw lines in those modules & 550,551 & \texttt{closure.py} \\
external imports: Mathlib modules / core modules & 50 / 1 & \texttt{closure.py} \\
\texttt{.lean} files under the family / outside the closure & 3,977 / 886 & this script \\
\texttt{axiom}, \texttt{sorry}, \texttt{admit}, \texttt{native\_decide}, \texttt{ofReduceBool} & 0, 0, 0, 0, 0 & \texttt{closure.py} \\
\texttt{implemented\_by}, \texttt{extern}, \texttt{unsafe}, \texttt{partial} & 0, 0, 0, 0 & \texttt{closure.py} \\
\texttt{elab}, \texttt{macro}, \texttt{syntax}, \texttt{simproc} & 0, 0, 0, 0 & \texttt{closure.py} \\
\texttt{set\_option} (both \texttt{quotPrecheck false}) & 2 & \texttt{closure.py} \\
\texttt{opaque} (all with bodies) & 9 (9) & \texttt{closure.py}, this script \\
\texttt{noncomputable} & 1,048 & \texttt{closure.py} \\
explicit \texttt{Classical.choice} / \texttt{propext} / \texttt{open Classical} & 4 / 4 / 8 & \texttt{closure.py} \\
\texttt{instance} / \texttt{local instance} & 148 / 40 & \texttt{closure.py} \\
\texttt{termination\_by} / \texttt{decreasing\_by} & 9 / 5 & \texttt{closure.py} \\
statement definitions compared / identical / identical up to \texttt{noncomputable} & 35 / 33 / 2 & \texttt{defdiff.py}, this script \\
theorem statement, characters (normalised / raw), identical & 1,124 / 1,353 & this script \\
\texttt{\#print axioms} commands: closure files / lines; family files / lines & 8 / 19; 17 / 35 & this script \\
recorded \texttt{depends on axioms} outputs in the clone & 0 & this script \\
challenge JSONs: total / \texttt{enable\_nanoda} false / true / absent & 418 / 413 / 2 / 3 & this script \\
CI workflow files (\texttt{lock-prs.yml} only) & 1 & this script \\
manifest packages / patch files / patches for Mathlib & 42 / 23 / 0 & this script \\
\bottomrule
\end{tabular}

\end{table}

\subsection{The statement is the challenge's}
\label{sec:proof-match}

The solution file does not import the challenge file:
\clonefile{SqrtLogSpace/Final.lean} begins (lines 1--3) with
\texttt{NativeLargeEntry}, \texttt{SqrtLogSpace.Statement} and
\texttt{SqrtLog.StatementCanonical}. The library re-declares every
definition the statement uses under the same fully qualified names,
spread over seven files: the SAT vocabulary in \clonefile{SAT.lean}, the
bit arithmetic in \clonefile{BinaryArithmetic.lean}, the nine encoding
functions in \clonefile{Encoding.lean}, the machine in
\clonefile{Computation.lean}, \texttt{runtimeScale} and
\texttt{runtimeBound} in \clonefile{SqrtLog/StatementScale.lean}, the
space instrumentation in \clonefile{SqrtLogSpace/VisitedTape.lean} and
\texttt{countWord} in \clonefile{SqrtLogSpace/EncodingBounds.lean}
(line~17, the challenge's line~147). The definition diff compares
\textbf{35} declarations: \textbf{33} are identical character for
character after whitespace normalisation, and the remaining two,
\texttt{runtimeScale} and \texttt{runtimeBound}, differ only in that the
challenge writes \texttt{noncomputable def} where the library opens a
\texttt{noncomputable section} (\clonefile{SqrtLog/StatementScale.lean},
line~15)---a difference the kernel does not see. The theorem statement
itself, from \texttt{theorem exact\_count\_and\_sat} to \texttt{:= by},
is \textbf{identical} in the two files, 1,353 characters raw and 1,124
after normalisation, at line~194 of the challenge and line~74 of
\clonefile{SqrtLogSpace/Final.lean}; the challenge's proof is the single
word \texttt{sorry} at its line~218.

This textual match is a sanity check and not the guarantee. The
comparator \citep{LeanComparator} compares the two statements at the
level of exported constants, with every declaration the statement
mentions matched, so a library definition that differed from the
challenge's under the same name would make the comparator fail rather
than pass silently. What the text confirms is that there is no such
difference to find, and that the shape Section~\ref{sec:statement} read
is the shape the library proves.

\subsection{The spine of the proof}
\label{sec:proof-spine}

\paragraph{The machine.} The theorem's machine is
\texttt{GuardedCompiler.primitive program} with the count tape
\texttt{256} and the answer tape \texttt{1216} (both \texttt{: Fin 1220};
\clonefile{SqrtLogSpace/Final.lean}, line~100), where \texttt{program}
is \texttt{Native.countingProgram cutoff k} \texttt{(NativeLeafProgram.command k)}
(line~40), a \texttt{BitProgram.Program 80 1220 labels} (line~11): a
first-order program over 80 one-bit registers, 1,220 Boolean stacks and
a finite label set, with twelve instruction kinds
(\clonefile{BitProgram.lean}, lines 10--26). \texttt{primitive}
(\clonefile{SqrtLogSpace/GuardedCompiler.lean}, line~131) is a
\texttt{noncomputable abbrev}: it takes the transition table
\texttt{transition p} (lines 30--47, one case per instruction) and the
stack selector \texttt{selected p} (lines 14--17), feeds them to the
stack compiler's controller, and renames the finite state type to
\texttt{Fin stateCount} through \texttt{Fintype.equivFin} (line~132).
The machine is thus a concrete table defined by pattern matching on the
program, composed with a classical renaming of states; the theorem is
existential, so the noncomputability costs it nothing, and it is the
reason no running time can be read off the statement
(Section~\ref{sec:statement}).

\paragraph{The time bound, and the constants.} For \texttt{cutoff < n},
\path{NativeLargeEntry.run} (\clonefile{SqrtLogSpace/NativeLargeEntry.lean},
lines 11--33) supplies a trace of length at most
\texttt{runtimeBound C A d L n} (line~27) with $1 \le C \wedge 1 \le A$
(line~12). Its exponent constant is \texttt{256 + At} (line~51), where
\texttt{At} comes from the leaf program's uniform contract
(\path{NativeLeafContract.uniform_contract}, line~39; the contract
gives $1 \le \mathtt{At}$, \clonefile{SqrtLogSpace/NativeLeafContract.lean},
line~793), and its cutoff is \texttt{max 20 (max Nleaf Ns)} (line~50),
with \texttt{Ns} obtained non-constructively from an eventual inequality
(\path{Filter.eventually_atTop}, line~48). The $2^{A\cdot\mathtt{runtimeScale}\,n}$
factor is assembled by \path{runtime_of_tree_and_leaf} (line~46)
from the DFS visit bound of \path{Native.outer_from_setup}
(\clonefile{SqrtLogSpace/NativeOuter.lean}, lines 53--56: visits at most
a square polynomial in $n$ times $2^{n/\mathtt{sparsifySchedule}\,n}$)
and the leaf time. For $n \le \mathtt{cutoff}$,
\path{Native.entry_small_primitive}
(\clonefile{SqrtLogSpace/NativeEntry.lean}, lines 81--100) gives
$\mathtt{time} \le 28700\,(L+1)^2$ (line~92) through the finite scanner
of \clonefile{SqrtLogSpace/FiniteFallback.lean}. The two branches are
folded by \texttt{runtime\_with\_fallback} and
\texttt{space\_with\_fallback}
(\clonefile{SqrtLogSpace/NativeOuterBudgets.lean}, lines 297 and 317)
into one triple, and the witnesses the final theorem exhibits are
\texttt{C + 28700, A + 1, S + 5022741, d + 2, e + 2}
(\clonefile{SqrtLogSpace/Final.lean}, line~41). So the proof as written
has an exponent constant of at least $258$ in front of $\rhon$, a
polynomial degree \texttt{d} that depends on $k$ (line~51: $\mathtt{db} + (4(k-2)+2) + \mathtt{df} + 2$) and a leading constant of at least
$28{,}701$; none of these is in the statement, which asserts only that
they are positive, and none of \texttt{At}, \texttt{db}, \texttt{df} is
computed anywhere. This confirms, from the proof side, what
Section~\ref{sec:statement} said from the statement side: the constants
are unknown, and the \emph{smallest} they could be is already far from
$1$.

\paragraph{Where a hidden cost model would live, and why it cannot.}
A proof of a running-time bound for a machine it compiles from a
higher-level program could, in principle, prove the bound against a cost
function of its own and leave the kernel to check only that function's
arithmetic. That is not what happens here. The compile lemmas conclude in
the comparator's own vocabulary: \path{GuardedCompiler.primitive_execution}
(\clonefile{SqrtLogSpace/GuardedCompiler.lean}, lines 158--162) turns a
source run of \texttt{n} steps into \texttt{RunsIn (primitive p).step
(7 * n)}, with the overhead \texttt{7} proved as a theorem
(\texttt{compiler\_overhead}, lines 152--153), and
\path{CompileSpace.compile_word_answer_space}
(\clonefile{SqrtLogSpace/CompileSpace.lean}, lines 125--143) produces a
\texttt{Computation.Trace (locatedStep (primitive p)) time \ldots} with
$\mathtt{time} \le 7\cdot\mathtt{sourceSteps}$ and
$\mathtt{visitedSpace}\ \mathtt{intermediate} \le s\,(\mathtt{bound} + \mathtt{input.length} + 15)$ on
every prefix (lines 134, 135 and 143); \path{primitive_located_step}
(\clonefile{SqrtLogSpace/GuardedCompiler.lean}, line~600) is the per-step
form. \texttt{Trace}, \texttt{locatedStep} and \texttt{visitedSpace} are
the challenge's definitions (its lines 87--90 and 149--189), and
\texttt{time} in the final theorem is literally the length of a
\texttt{Computation.Trace} of \texttt{PrimitiveMachine.step} (its
line~201). No cost function of the compiler's appears in the statement;
a wrong cost model would make these lemmas unprovable, not the theorem
false. This is the structural point on which the rest of the paper
leans: whatever the proof's internal accounting, the number the theorem
bounds is the number of steps of the comparator's machine.

\paragraph{Correctness of the count.} \path{Native.outer_from_setup}
(\clonefile{SqrtLogSpace/NativeOuter.lean}, lines 50--51) yields
\texttt{final.stacks ABI.countStack = countWord (SAT.count F)} and
\texttt{final.stacks ABI.answerStack = [decide (0 < SAT.count F)]} from
the canonical DFS execution, and
\path{Native.counting_primitive_from_entry} (line~73) transports
both to the primitive machine's tapes; \texttt{countWord\_toNat} and
\texttt{countWord\_count\_length} supply the decoding and length clauses
of the statement. These are theorems about the concrete step function.
The arithmetic that makes them true---sparsification, the ring and
moment computations, the DFS---is the content of the other three
thousand files, and the only claim made about it here is the one the
kernel makes: it elaborates (Section~\ref{sec:replay}).

\subsection{Toolchain, Mathlib, and the build hooks}
\label{sec:proof-build}

The toolchain is fixed (\clonefile{lean/lean-toolchain}:
\texttt{leanprover/lean4:v4.34.1}; \clonefile{lean/lakefile.lean}
line~8: \texttt{fixedToolchain := true}). The manifest lists 42
packages; Mathlib is pinned to upstream revision
\texttt{d13f23b723b8a846827a245b89c10fc7d3f11612} at
\path{leanprover-community/mathlib4}, and the audit of 2026-11-05
confirmed that revision on GitHub as the upstream commit that bumps the
toolchain to \texttt{v4.34.1}. The repository carries 23 compatibility
patches under \clonefile{lean/patches/}, \textbf{none for Mathlib};
every one is for a third-party package (\texttt{AINTLIB},
\texttt{PrimeNumberTheoremAnd}, \texttt{carleson}, \texttt{iut}, \ldots) that
this closure does not import. The lakefile has two hooks that run
arbitrary IO when Lake configures or updates the package: a
\texttt{run\_cmd} block (lines 290--304) that evaluates dependency
declarations with \texttt{unsafe env.evalConstCheck} and calls
\path{preparePatchedDependencies}, which clones and patches eleven
dependencies, and a \texttt{post\_update} hook (line~306 onward) that
verifies or applies the patches for the rest. They touch only
\texttt{.lake/packages}, and only non-Mathlib packages; they are one
reason the inspection did not build, and the reason the replay of
Section~\ref{sec:replay} runs inside a sandbox with the network on.

\subsection{What the repository does not record}
\label{sec:proof-missing}

Three things one would expect beside a verified refutation of a
forty-year-old hypothesis are absent, and their absence is a fact about
evidence rather than about the proof.

\begin{itemize}
\item \textbf{No axiom audit.} \clonefile{SqrtLogSpace/Final.lean} ends
  with five \texttt{\#print axioms} commands (lines 197--201), and seven
  other files of the closure carry fourteen more (nineteen commands in
  eight files; the family's directory holds seventeen such files, but
  nine of them are among the 886 never built). Their output exists only
  while the build runs and is recorded nowhere: a search of the whole
  clone for the string \texttt{depends on axioms} finds nothing.
\item \textbf{No continuous integration.} \clonefile{.github/workflows/}
  holds one file, \texttt{lock-prs.yml}, which locks pull-request
  conversations. Nothing builds or checks the Lean.
\item \textbf{No second kernel for this family.} The comparator's JSON
  (\clonefile{lean/ComparatorChallenges/SubexponentialSAT.json}) names the
  one theorem, an empty \texttt{definition\_names}, the three permitted
  axioms \texttt{propext}, \texttt{Quot.sound}, \texttt{Classical.choice},
  and \texttt{enable\_nanoda: false}; the registry
  \clonefile{lean/formalization.yaml} (lines 1683--1685) points at the
  same declaration and file. Of the collection's 418 challenge JSONs,
  413 set \texttt{enable\_nanoda} to \texttt{false}, two to
  \texttt{true}, and three omit the key. With it off, the collection's own
  check of this family rests on one kernel implementation. The
  comparator's README, as fetched by the audit, describes the check as:
  build challenge and solution in a sandbox, export both with
  \texttt{lean4export}, match each named theorem's statement with every
  declaration it uses, confine the proof to the permitted axioms, and
  have the Lean kernel accept it; with nanoda on, the exported
  environment is replayed through an independent kernel in Rust as well
  \citep{Nanoda}. Section~\ref{sec:related} records the versions.
\end{itemize}

Two smaller consistency facts, for the record. The manuscript's source is
dated ``17 October 2026, revised 25 October 2026''
(\clonefile{preprints/The-Exponential-Time-Hypothesis-is-False-October-25-2026/build/main_full_version.tex}, line~25); the
collection's history file mentions the family once, in the sentence that
names the working group at whose request the materials were released
(\clonefile{history.md}, line~18), and its only ``Formalization Fix''
entry (line~7) concerns another family. And the scope note
\clonefile{lean/docs/378.md} states the verified exponent as
$n(\log\log n)^2/\sqrt{\log n}$ (line~9), which is what the library
proves and what Section~\ref{sec:beyond} compares with the manuscript's.

\subsection{The trust boundary, joined with the replay}
\label{sec:proof-trust}

The audit's verdict listed what the proof's trust rested on before
anything was run: (i) OpenAI's unpublished build, since nothing showed
the closure type-checks; (ii) the absence of any recorded axiom output;
(iii) the Lean kernel of \texttt{v4.34.1} and Mathlib \texttt{d13f23b7}
as upstream; (iv) the comparator run, not reproduced; (v) no second
kernel for this family. Table~\ref{tab:proof-trust} says which of the
five Section~\ref{sec:replay} moved.

\begin{table}[t]
\centering\small
\caption{The trust boundary before and after the replay of
Section~\ref{sec:replay}.}
\label{tab:proof-trust}
\begin{tabular}{@{}p{0.30\linewidth}p{0.64\linewidth}@{}}
\toprule
item & status after the replay \\
\midrule
(i) the build & \textbf{moved}: the 3,091-module closure built from source in our sandbox at the pinned toolchain, exit~0, no \texttt{error}, no \texttt{sorry}. \\
(ii) the axiom output & \textbf{moved}: \texttt{\#print axioms} of the comparator's theorem printed \texttt{[propext, Classical.choice, Quot.sound]}, and the nineteen in-build axiom lines say the same; the output is on record. \\
(iii) kernel and Mathlib & \textbf{moved in part}: Mathlib was fetched from upstream at the pinned revision and not patched, so ``as upstream'' is now observed rather than assumed; the Lean kernel of \texttt{v4.34.1} is still trusted as upstream, now beside a second kernel (v). \\
(iv) the comparator run & \textbf{not moved}: the statement-level match through \texttt{lean4export} was not performed; the match this paper relies on is the textual one of Section~\ref{sec:proof-match}. \\
(v) a second kernel & \textbf{moved}: \texttt{nanoda\_lib} checked the exported dependencies of the theorem, 95,176 declarations, with no errors and the three axioms; what this adds and what it does not is stated in Section~\ref{sec:replay}. \\
\bottomrule
\end{tabular}
\end{table}

After the replay, then, the proof's trust rests on: the Lean kernel of
\texttt{v4.34.1} \emph{and} nanoda agreeing on the exported dependencies
of the theorem; \texttt{lean4export} exporting what the environment
holds; upstream Mathlib at the pinned revision; and the textual, not
export-level, match between the solution's statement and the
challenge's. Everything the inspection could find by reading, the build
and the two kernels confirmed; nothing the inspection could not settle
by reading has been settled by anything else.

\subsection{What this does not claim}
\label{sec:proof-not}

\begin{itemize}
\item \textbf{The arithmetic was not read.} The inspection followed the
  spine from \clonefile{SqrtLogSpace/Final.lean} down to the compile
  lemmas and swept the whole closure by keyword; the sparsification, ring
  and DFS content of the other three thousand files was not read. The
  argument above is that its correctness is kernel-enforced, not that we
  verified it.
\item \textbf{Grep is necessary, not sufficient.} Zero \texttt{sorry}
  and zero \texttt{axiom} by text search cannot show that a file
  elaborates; only the build does (Section~\ref{sec:replay}). Conversely
  a keyword count is over source text at one commit, and a declaration
  can be unsound only through the mechanisms searched for; we searched
  for the ones we know.
\item \textbf{The textual statement match is not the comparator's.} The
  comparator matches exported constants; we matched source text after
  whitespace normalisation and found it identical. A difference the
  normalisation hides would have to be invisible to the kernel as well,
  but that is an argument, not a run.
\item \textbf{Nothing about the constants' values.} The bounds
  $A \ge 258$ and $C \ge 28{,}701$ are what the proof's witnesses are
  built from; they are lower bounds on what this proof exhibits, not
  statements about any machine's running time, and a different proof of
  the same statement could exhibit different ones.
\item \textbf{Nothing about the Lean kernel, Mathlib or the toolchain.}
  All three are trusted as upstream at the pinned versions; the second
  kernel of Section~\ref{sec:replay} is evidence about the first, not a
  proof of either.
\item \textbf{Nothing about what the theorem says.} That is
  Section~\ref{sec:statement}; a sound proof of a weak bound is a sound
  proof of a weak bound.
\end{itemize}
\par\endgroup
